% ৩.১৩ সত্যক সারণি (Truth Table)
\section{Truth Table}

\subsection{সংজ্ঞা}
\begin{itemize}
\item Input combination অনুযায়ী output নির্ণয়
\end{itemize}

\subsection{প্রশিক্ষণ}
\begin{itemize}
\item ১,২,৩ input truth table উদাহরণ
\item Logic gate যাচাই
\end{itemize}

\begin{learningbox}
এই টপিক শেষে শিক্ষার্থী পারবে:
\begin{itemize}
\item input সংখ্যা অনুযায়ী truth table row নির্ণয় করতে।
\item Boolean expression থেকে output column পূরণ করতে।
\item logic gate-এর কাজ truth table দিয়ে যাচাই করতে।
\end{itemize}
\end{learningbox}

\begin{definitionbox}{Truth table}
কোনো logic expression বা circuit-এর সম্ভাব্য সব input combination এবং প্রতিটির জন্য output দেখানো সারণিকে truth table বলা হয়।
\end{definitionbox}

\begin{keypointbox}{Row সংখ্যা}
\begin{itemize}
\item 1 input হলে row = $2^1=2$
\item 2 input হলে row = $2^2=4$
\item 3 input হলে row = $2^3=8$
\item সাধারণভাবে $n$ input হলে row = $2^n$
\end{itemize}
\end{keypointbox}

\begin{examplebox}{Expression: $Y=A+B\cdot C$}
প্রথমে $B\cdot C$ column তৈরি করো, তারপর $A$-এর সাথে OR করো। intermediate column রাখলে ভুল কম হয়।
\end{examplebox}

\begin{center}
\fbox{\parbox{0.88\textwidth}{
\centering
\textbf{Truth table placeholder}\\[4pt]
Columns: $A$, $B$, $C$, $B\cdot C$, $Y=A+B\cdot C$\\
৮টি row পূরণ করার জায়গা।
}}
\end{center}
